% !TeX root = se100_the_ontological_neutrality_theorem.tex
\section{Conclusion}
\label{sec:conclusion}

This paper has examined the structural conditions under which
an ontology can function as a neutral substrate
across admissible interpretive frameworks under persistent disagreement.
The central problem is not semantic heterogeneity alone,
nor the absence of a shared vocabulary.
It is the need to maintain stable shared reference
while allowing legal, causal, analytic, political,
or normative frameworks to reach incompatible conclusions
about the same referents.

The result established here is conditional but strong.
Under the neutrality requirements of interpretive non-commitment
and extension stability,
under the contestability assumption for causal and normative propositions,
and under the scope condition
that the substrate's
framework-invariant referential structure, including its
referential regime,
is invariant
across the admissible frameworks under consideration,
a substrate ontology satisfies neutrality if and only if
its foundational layer is restricted to framework-invariant
referential structure and
does not assert
causal or normative conclusions as substrate-layer facts.

This theorem identifies a categorical boundary
in ontology design.
A neutral substrate is not empty,
contentless, or epistemically indifferent.
It must assert enough structure to support accountability:
entities, occurrences, institutional artifacts,
referential regimes, provenance, attribution,
and other framework-invariant structural relations.
In particular, it must provide the referential regimes
by which relevant referents are individuated,
co-referred, and tracked across time, records,
jurisdictions, or interpretive contexts.

At the same time,
the substrate must not assert causal or normative conclusions
as foundational facts.
Claims that one event caused another,
that an actor was responsible,
that an action was permitted or prohibited,
or that a system was compliant or noncompliant
belong to interpretive frameworks.
They may be represented by the substrate as reified,
attributed, provenance-bearing claims,
but their truth must not be asserted at the substrate layer.
The exclusion is therefore at the level of assertion,
not vocabulary.

This distinction is the core design constraint.
A substrate may record that a framework made a causal or normative claim;
it may record who made the claim,
when it was made,
what referents it concerns,
and what provenance supports it.
But the substrate may not convert that claim
into a substrate-layer ontological fact
without forfeiting neutrality.
Representing discourse is compatible with neutrality.
Endorsing the contested content of that discourse
at the foundational layer is not.
The theorem also clarifies the scope of neutrality.
Neutrality is not global.
A substrate is not neutral with respect to
the referential regimes it asserts.
It must take substantive positions on individuation,
co-reference, and persistence
where those positions are required for stable accountability.
If admissible frameworks disagree about those referential conditions,
then neutrality in the sense defined here is unavailable
until the referential contestation is resolved or scoped out.
This limitation is not a defect of the theorem;
it is a boundary condition on the possibility of a neutral substrate.

The resulting architecture separates two tasks
that accountability systems often conflate.
The substrate provides stable referential ground:
what was recorded,
what it concerns,
who asserted it,
and what provenance accompanies it.
Interpretive frameworks provide explanation,
evaluation, judgment, responsibility assignment,
and causal or normative conclusion.
This separation allows competing frameworks to reason
over the same substrate without forcing the substrate itself
to adjudicate among them.

The practical implication is that neutral substrates
should be designed as pre-causal and pre-normative foundations.
Pre-causal does not mean acausal,
and pre-normative does not mean normatively irrelevant.
Rather, causal and normative reasoning are externalized
to framework layers where their assumptions are explicit,
attributed, contestable, and revisable.
This preserves accountability under disagreement
by making competing interpretations representable
without destabilizing the shared record.

The contribution of the paper is therefore not a new ontology,
nor a protocol for resolving disagreement.
It is a structural constraint on any ontology intended to function
as a neutral substrate in accountability contexts.
Where persistent interpretive disagreement is expected,
and where the substrate must remain stable across admissible frameworks,
causal and normative conclusions cannot be asserted as substrate-layer facts.
They must be externalized as framework-level conclusions
or represented as reified claims with provenance.

In this sense,
ontological neutrality is not a lack of commitments.
It is a disciplined allocation of commitments across layers.
The substrate commits to stable reference.
Frameworks commit to interpretation.
Accountability is preserved not by eliminating disagreement,
but by giving disagreement a stable shared object
over which it can be expressed, evaluated, and contested.

\medskip
\nopagebreak
\section*{Acknowledgements}

The author gratefully acknowledges the philosophical foundations,
research communities, and systems engineers whose prior work and engagement
informed the development of the ideas presented here.

Portions of this work were developed
using computer-assisted tools during manuscript preparation.
Generative language models were used to assist with editing, formatting,
and consistency checking.
All conceptual framing, formal development, results, interpretations, and conclusions
are the author's own.

The author reviewed all suggestions and
takes full responsibility for the content of this work.
